\documentclass[../../../include/open-logic-section]{subfiles}

\begin{document}

\olfileid{sth}{choice}{tarskiscott}
\olsection{ترفند تارسکی--اسکات}

در \olref[cardinals][cardsasords]{defcardinalasordinal}، کاردینال‌ها را
به‌صورت اعداد ترتیبی تعریف کردیم. برای این کار اصل خوش‌ترتیبی را فرض
کردیم. تنها دلیل این انتخاب آن بود که این تعریف «استاندارد رایج» است.

پس پیش از پرداختن به هر یک از مسائل فلسفی پیرامون خوش‌ترتیبی، باید روشن
کنیم که \emph{می‌توانیم} از استاندارد رایج فاصله بگیریم و نظریه‌ای برای
کاردینال‌ها \emph{بدون} فرض خوش‌ترتیبی بپرورانیم. همچنان می‌توانیم
تعریف‌های $\cardeq{A}{B}$، $\cardle{A}{B}$ و $\cardless{A}{B}$ را،
همان‌گونه که در \olref[sfr][siz][]{chap} آمده‌اند، به کار ببریم. فقط به
مفهوم تازه‌ای از \emph{کاردینال} نیاز خواهیم داشت.

اندیشه‌ای ساده‌انگارانه این است که بکوشیم کاردینالیتهٔ $A$ را چنین تعریف
کنیم:
\[
	\Setabs{x}{\cardeq{A}{x}}.
\]
این تعریف را می‌توان با تعریف فرگه از $\# x Fx$، که در پایان
\olref[cardinals][hp]{sec} طرح شد، مقایسه کرد. به دلایلی که در آن‌جا به
آن‌ها اشاره کردیم، این تعریف شکست می‌خورد. برای نمونه، بگذارید
$A=\{\emptyset\}$. هر مجموعهٔ تک‌عضوی با $A$ هم‌توان است، اما در هر
مرحلهٔ جانشینِ سلسله‌مراتب مجموعه‌های تک‌عضوی تازه‌ای پدید می‌آیند
(کافی است مجموعهٔ تک‌عضویِ مرحلهٔ پیشین را در نظر بگیرید). بنابراین
$\Setabs{x}{\cardeq{A}{x}}$ وجود ندارد، زیرا نمی‌تواند رتبه‌ای داشته
باشد.

برای رفع این مشکل، از ترفندی منسوب به تارسکی و اسکات استفاده
می‌کنیم:\footnote{یادآوری: همهٔ فرمول‌ها می‌توانند پارامتر داشته باشند،
مگر آن‌که خلاف آن صریحاً گفته شود.}

\begin{defn}[Tarski--Scott]
برای هر فرمول $\phi(x)$، بگذارید $[ x : \phi(x)] $ مجموعهٔ همهٔ
$x$هایی با کمترین رتبهٔ ممکن باشد که $\phi(x)$ را ارضا می‌کنند (و اگر
هیچ $x$ای $\phi$ را ارضا نکند، بگذارید برابر $\emptyset$ باشد).
\end{defn}

باید بررسی کنیم که این تعریف موجه است. در $\ZF$،
\olref[spine][foundation]{zfentailsregularity} تضمین می‌کند که
$\setrank{x}$ برای هر $x$ وجود دارد. حال اگر موجوداتی $\phi$ را ارضا
کنند، کمترین $\alpha$ را که
$(\exists x\subseteq V_\alpha)\phi(x)$ برای آن برقرار است برمی‌گزینیم.
بنا بر کمینگی $\alpha$، برای هر رتبهٔ پایین‌تر هیچ مصداقی از $\phi$
وجود ندارد؛ یعنی
$(\forall \beta \in \alpha)(\forall x \subseteq V_\beta)
\lnot \phi(x)$. سپس می‌توانیم با اصل جداسازی،
$[x : \phi(x)]$ را به‌صورت
$\Setabs{x \in V_{\alpha+1}}{\phi(x)}$ تعریف کنیم.

اکنون که ترفند تارسکی--اسکات را توجیه کرده‌ایم، می‌توانیم آن را برای
تعریف مفهومی از کاردینالیته به کار گیریم:

\begin{defn}
کاردینالیتهٔ \textsc{ts} مجموعهٔ $A$ عبارت است از
$\text{tsc}(A) = [x : \cardeq{A}{x}]$.
\end{defn}

تعریف کاردینال \textsc{ts} از خوش‌ترتیبی استفاده نمی‌کند. اما حتی بدون
آن اصل نیز می‌توانیم نشان دهیم که \emph{کاردینال‌های \textsc{ts}} کمابیش
مانند \emph{کاردینال‌ها}ی تعریف‌شده در
\olref[cardinals][cardsasords]{defcardinalasordinal} رفتار می‌کنند. برای
مثال، اگر \olref[cardinals][cardsasords]{lem:CardinalsBehaveRight} و
\olref[card-arithmetic][opps]{lem:SizePowerset2Exp} را با زبان
کاردینال‌های \textsc{ts} بازبیان کنیم، اثبات‌ها در $\ZF$ و بدون فرض
خوش‌ترتیبی به همان صورت پیش می‌روند.

در همین زمینه شایان ذکر است که با ترفند تارسکی--اسکات می‌توان نظریه‌ای
برای اعداد ترتیبی نیز پروراند. اگر $\tuple{A, <}$ یک خوش‌ترتیب باشد،
بگذارید $\text{tso}(A, <) = [\tuple{X, R} :
\ordeq{\tuple{A, <}}{\tuple{X, R}}]$. برای شرح بیشتر این شیوهٔ پرداختن
به کاردینال‌ها و اعداد ترتیبی، بنگرید به
\citet[chs.~9--12]{Potter2004}.

\end{document}
